\documentclass[11pt,UTF8,fontset=fandol]{ctexart}
\usepackage{amsmath,amssymb,amsfonts}
\usepackage{graphicx,float,caption,subcaption}
\usepackage{booktabs,array,multirow,makecell}
\usepackage{geometry,enumitem,xcolor}
\usepackage[hidelinks]{hyperref}
\usepackage{url}
\geometry{margin=2.3cm}
\graphicspath{{/Users/yye/project/论文/assets/retool/}}
\captionsetup{font=small,labelfont=bf}
\setlength{\parskip}{0.4em}

% 自定义小白知识卡盒子
\usepackage[most]{tcolorbox}
\newtcolorbox{infobox}[1]{colback=blue!4!white,colframe=blue!45!black,
  fonttitle=\bfseries,title={#1},boxrule=0.6pt,arc=2pt,left=4pt,right=4pt,top=3pt,bottom=3pt}
\newtcolorbox{takeaway}{colback=green!4!white,colframe=green!40!black,
  boxrule=0.6pt,arc=2pt,left=5pt,right=5pt,top=4pt,bottom=4pt}

\title{\textbf{ReTool：用强化学习教大模型“会用工具”——\\在推理中策略性地调用代码解释器}}
\author{Jiazhan Feng, Shijue Huang, Xingwei Qu, Ge Zhang, Yujia Qin,\\
Baoquan Zhong, Chengquan Jiang, Jinxin Chi, Wanjun Zhong\\
\textbf{ByteDance Seed}（字节跳动 Seed 团队）}
\date{}

\begin{document}
\maketitle

\begin{center}
\fbox{\parbox{0.92\linewidth}{\small
\textbf{英文原标题：} ReTool: Reinforcement Learning for Strategic Tool Use in LLMs\\
\textbf{发表：} arXiv:2504.11536v2，2025 年 4 月 \quad|\quad \textbf{引用数：} 297（Semantic Scholar, 2026-06-07）\\
\textbf{代码 / 项目主页：} \url{https://github.com/ReTool-RL/ReTool} \quad \url{https://retool-rl.github.io/}\\[2pt]
\textbf{一句话定位：} 这篇论文要解决的是——\emph{“会推理”的大模型为什么算不准数、解不对复杂方程}。它的答案是：
用强化学习（RL）让模型\emph{自己学会}在文字推理过程中“随手写一段代码 → 真跑一下 → 把运行结果接着往下想”，
而“什么时候该写代码、写什么代码”全靠 RL 摸索出来，不靠人写死规则。
}}
\end{center}

\vspace{0.5em}

%========================= 小白知识卡 =========================
\section{读这篇前要懂的几个概念（小白知识卡）}

下面这些是读懂这篇论文的“前置知识”。如果你是 AI 应用开发的新手，先花两分钟扫一遍。

\begin{infobox}{1. 大语言模型（LLM）与思维链（CoT）}
大语言模型（large language model, LLM）就是 GPT、DeepSeek 这类会“说话”的模型。
\textbf{思维链}（chain-of-thought, CoT）指的是：让模型在给最终答案前，先用自然语言一步步“把思路写出来”，
就像做数学题写草稿。写草稿往往能让模型答得更准。本文研究的“推理模型”（如 OpenAI o1、DeepSeek R1）就是
特别擅长写长草稿、还能自我检查的模型。
\end{infobox}

\begin{infobox}{2. 强化学习（RL）与 PPO}
\textbf{强化学习}（reinforcement learning, RL）是一种“通过试错+奖惩”来训练模型的范式：模型先生成一批答案
（这一步叫\textbf{rollout}，可以理解为“让模型自由发挥跑一遍”），然后根据答案好坏给一个\textbf{奖励}（reward），
好的行为被强化、坏的被抑制。\textbf{PPO}（Proximal Policy Optimization，近端策略优化）是目前最常用的 RL 算法之一，
它的核心技巧是：每次更新不让模型“变化太大”（用一个 clip 裁剪机制限制步子），从而训练更稳。
\end{infobox}

\begin{infobox}{3. RLVR：用“可验证的结果”当奖励}
本文用的是\textbf{结果奖励}（outcome reward）：不管模型中间过程怎么折腾，\emph{只看最终答案对不对}，对了给
$+1$、错了给 $-1$。因为数学题答案唯一、可以程序自动判分，所以这类奖励叫“可验证奖励”（这一思路常被称为
RLVR, RL with Verifiable Rewards）。好处是\textbf{不需要人工标注}每一步，也不容易被模型“钻空子刷分”（reward hacking）。
\end{infobox}

\begin{infobox}{4. 工具集成推理（TIR）与代码解释器（CI）}
\textbf{代码解释器}（code interpreter, CI）就是一个能真正运行 Python 代码的“沙箱”（sandbox，隔离的安全运行环境）。
\textbf{工具集成推理}（tool-integrated reasoning, TIR）指的是：让模型在推理时不要光靠脑子硬算，而是把
“枚举、验证、精确计算”这种活儿交给代码去做。本文就是把代码解释器接进了模型的推理回路里。
相关的早期思路还有 PoT（Program-of-Thought，程序化思维）、PAL（程序辅助语言模型）。
\end{infobox}

\begin{infobox}{5. 冷启动 SFT（cold-start supervised fine-tuning）}
RL 训练前，模型最好已经“大致会”做这件事，否则一上来全是随机乱试、学不动。所以本文先用一批
高质量示范数据做\textbf{监督微调}（supervised fine-tuning, SFT），教模型“代码该长什么样、什么时候插进来”，
给后续 RL 一个好的起点——这一步就叫\textbf{冷启动}（cold-start）。
\end{infobox}

%========================= 摘要全译 =========================
\section{摘要（全文翻译）}

用强化学习（RL）训练的推理模型（例如 DeepSeek R1）在\emph{文字推理}上表现出色，但在需要\emph{结构化求解}的场景里
却力不从心——比如几何推理、需要简洁计算的题目、或复杂方程求解；而这些恰恰是代码解释器（code interpreter, CI）
这类计算工具的强项。为弥合这一鸿沟，我们提出 \textbf{ReTool}，用“工具集成学习”增强长链推理，它有两个关键特性：
（1）在自然语言推理过程中\textbf{动态地交错插入实时代码执行}；（2）一套\textbf{自动化的 RL 范式}，允许策略在 rollout 中
进行多轮、实时的代码执行，并\emph{仅凭结果反馈}教模型学会“何时、如何调用工具”。

ReTool 采用一套系统化的训练框架：先用合成的冷启动数据，生成“带代码的长链推理轨迹”，用来微调基座模型；
随后的 RL 训练以\emph{任务结果作为奖励}，迭代地打磨模型的工具使用策略，使其能够\textbf{自主发现最优的工具调用模式，
而无需人类先验}。在具有挑战性的数学奥赛基准 AIME 上的实验证明了 ReTool 的优越性：我们的 32B 模型仅用
\textbf{400 个训练步}就在 AIME2024 上达到 \textbf{67\% 的准确率}，无论在效率还是性能上都超过了文字 RL 基线
（40\% 准确率，1080 步）。更值得注意的是，在扩展设置下 ReTool-32B 取得了 \textbf{72.5\% 的准确率，比 OpenAI 的
o1-preview 高出 27.9\%}。进一步分析揭示出一些\textbf{涌现行为}（emergent behaviors），例如代码自我纠错，标志着模型
自主掌握自适应工具使用的“顿悟时刻”（aha moment）。这些发现凸显了“结果驱动的工具集成”在推进复杂数学推理上的
潜力，也为混合的“神经—符号”系统（neuro-symbolic systems）提供了新见解。

\smallskip
\noindent\textit{日期：2025 年 4 月 15 日 \quad 项目主页：}\url{https://retool-rl.github.io/}

%========================= 逐节精读 =========================
\section{逐节精读}

\subsection{第 1 节 引言：为什么“会推理”还不够？}

作者先承认现状：RL 已经成为提升 LLM 推理能力的主流范式，o1、DeepSeek R1 这些模型靠它学会了自我纠错、
更审慎地思考，甚至显露出“元认知”（metacognition，即不仅会想，还会\emph{监督和修正}自己的思考）的苗头。

\textbf{但有个老大难：纯文字推理算不准数。} 几何题、需要精确数值计算的题、复杂方程，这些场景下模型光靠“语言模式”
硬算，特别容易出小错，而且错误会层层累积（论文叫 compounding error，复合误差）。相比之下，代码解释器提供的是一个
\emph{形式化、可执行}的接口，擅长枚举、验证、精确计算，还能让模型通过“写程序去探索”来扩大解题的搜索空间。

\textbf{那为什么不直接教模型用工具呢？} 之前的做法主要是\emph{提示工程}和\emph{监督微调}（SFT）。问题在于：SFT 只是在
\emph{模仿}人精心准备的那批数据，学到的是固定套路，一旦题目变样就不灵了，模型可能\emph{乱用工具}或退化成脆弱的
经验法则。作者的洞察是：\textbf{RL 才是正解}——它让模型自由探索推理轨迹，靠“结果好不好”这个反馈来学策略，
从而能学会一些精细的行为，比如“代码跑错了怎么从错误中恢复”“长推理里到底该在哪一步调用代码”。

\textbf{ReTool 的两大组件}（这是全文骨架，务必记住）：
\begin{enumerate}[leftmargin=2em,itemsep=2pt]
  \item \textbf{冷启动数据构造}：做一个高质量数据集，明确示范“何时、如何调用代码解释器”，让模型先具备基本的
  工具使用与“看懂运行结果”的能力。
  \item \textbf{工具增强的 RL}：在此基础上用 RL 训练，让模型\emph{自己发现}最优的工具操控策略，并通过结果奖励
  不断调整行为——这是 SFT 单独做不到的。在长链推理中，策略模型会灵活地写代码块、从沙箱式代码解释器拿到
  实时运行结果，再据此继续往下想。
\end{enumerate}

\textbf{引言里抢先剧透的结果：} 基于 Qwen2.5-32B-Instruct，ReTool 仅 400 步就在 AIME2024 上达 67.0\%，
而文字 RL 基线跑 1080 步才 40.0\%。换更强的基座 DeepSeek-R1-Distill-Qwen-32B 后还能更高，超过 QwQ-32B-Preview、
s1-32B、o1-preview 等强基线。另外，光是冷启动模型（还没 RL）就有 40.9\%，已经和文字 RL 基线打平、远超原始
Qwen2.5-32B-Instruct 的 26.7\%——说明\emph{数据集本身就有效地捕捉到了工具使用模式}。

\subsection{第 2 节 方法：ReTool 是怎么做的}

\subsubsection{2.1 总览：两个阶段}
整个方法分两步走：\textbf{(1) 冷启动 SFT}（给 RL 一个稳健的初始化）；\textbf{(2) 带“交错代码执行”的 RL}
（专门训练模型在推理中恰当地选择和使用工具）。

\subsubsection{2.2 冷启动：打好工具集成推理的底子}
做法是设计一条数据流水线：
\begin{itemize}[leftmargin=1.6em,itemsep=2pt]
  \item \textbf{先攒纯文字推理数据}：从多个来源（含开源的 OpenThoughts 等）收集数学推理数据，再用\emph{“人审 + DeepSeek-R1
  评估”双重验证}过滤掉无效样本，得到高质量文字数据集，记作 $D_{\text{init}}$。
  \item \textbf{再自动转成“带代码”的数据}：用一个结构化提示模板，把原推理过程里那些\emph{适合用代码做的手工计算步骤}，
  替换成对应的代码片段 + 解释器运行结果。然后做\emph{两阶段验证}：第一阶段查\textbf{格式}（保证语法一致、可读，便于后续 RL
  里识别“何时该触发工具调用”）；第二阶段查\textbf{答案}（最终答案不对的样本直接扔掉）。最后得到带代码的长链推理数据集
  $D_{\text{CI}}$。
  \item \textbf{在 $D_{\text{CI}}$ 上做 SFT}，让模型学会“何时、如何调用代码解释器”。
\end{itemize}
\emph{白话总结}：先用大模型把“人脑硬算的步骤”改写成“写代码算”，再严格筛掉算错的，拿这批干净示范喂给模型打底。

\subsubsection{2.3 ReTool 的 RL：策略性工具使用}

\paragraph{训练算法（PPO）。} ReTool 基于 PPO 算法，优化目标如下（公式 1）：
\begin{equation*}
J_{\text{PPO}}(\theta)=\mathbb{E}_{(q,a)\sim D,\,o_{\le t}\sim\pi_{\theta_{\text{old}}}(\cdot|q)}
\left[\min\!\left(\frac{\pi_\theta(o_t|q,o_{<t};\,\text{CI})}{\pi_{\theta_{\text{old}}}(o_t|q,o_{<t};\,\text{CI})}\hat{A}_t,\;
\operatorname{clip}\!\left(\frac{\pi_\theta(o_t|q,o_{<t};\,\text{CI})}{\pi_{\theta_{\text{old}}}(o_t|q,o_{<t};\,\text{CI})},1-\varepsilon,1+\varepsilon\right)\hat{A}_t\right)\right]
\end{equation*}

\textbf{逐符号拆解（看懂这一个公式就懂 PPO 了）：}
\begin{itemize}[leftmargin=1.6em,itemsep=1pt]
  \item $\pi_\theta$：当前正在训练的\textbf{策略模型}（policy model），$\theta$ 是它的参数；$\pi_{\theta_{\text{old}}}$ 是更新前的\textbf{旧策略}（参考模型）。
  \item $q$ 是问题，$a$ 是标准答案，$D$ 是数据集；$o_t$ 是模型在第 $t$ 步生成的 token，$o_{<t}$ 是它前面已生成的内容。
  \item \textbf{关键点 }$\boldsymbol{\pi_\theta(o_t|q,o_{<t};\text{CI})}$：这里多了一个 “$;\text{CI}$”，意思是\emph{生成过程中夹带了代码解释器的实时执行与反馈}——
  这正是 ReTool 对标准 PPO 的改造：rollout 不再是纯文字，而是“文字 + 真跑代码”。
  \item $\dfrac{\pi_\theta(\cdot)}{\pi_{\theta_{\text{old}}}(\cdot)}$ 是\textbf{重要性比率}（新旧策略对同一动作的概率之比）；$\hat{A}_t$ 是\textbf{优势函数}
  （advantage，衡量“这一步比平均水平好多少”）。
  \item $\operatorname{clip}(\cdot,1-\varepsilon,1+\varepsilon)$ 把这个比率\textbf{裁剪}在 $[1-\varepsilon,1+\varepsilon]$ 区间内，再取 $\min$。
  这就是 PPO 的“别一步迈太大”的安全阀，保证训练稳定。
\end{itemize}

\paragraph{奖励设计（公式 2）。} 为了教模型“何时、如何用工具”，作者用了一个极简的\emph{基于规则的准确率奖励}：
要求模型把最终答案放进指定格式（如 \verb|\boxed{}|）里，便于程序自动判分。奖励定义为：
\begin{equation*}
R(a,\hat{a})=
\begin{cases}
\;\;1, & \text{若 } \mathrm{is\_equivalent}(a,\hat{a})\ \text{（预测答案与标准答案等价）}\\[2pt]
-1, & \text{否则}
\end{cases}
\end{equation*}
其中 $a$ 是标准答案、$\hat{a}$ 是模型预测答案。\textbf{为什么这么简单？} 作者特意只看最终对错、\emph{不给“代码能否运行”单独加分}，
就是为了\textbf{避免 reward hacking（钻奖励空子）}、并鼓励模型探索更多样的解题行为。这也呼应了 RLVR 的思路：
奖励越简单越纯粹，越不容易被带偏。

\paragraph{带交错代码执行的 Rollout（方法的灵魂）。} 这是 ReTool 跟普通 RL 最不同的地方（见后文图~\ref{fig:arch}）。普通 rollout
只生成文字；ReTool 的 rollout 让策略模型和外部代码沙箱\emph{协作}，产出“文字 + 代码 + 实时解释器反馈”的混合内容。
具体流程是：
\begin{enumerate}[leftmargin=1.8em,itemsep=1pt]
  \item 模型用 \verb|<code></code>| 标签把要执行的代码圈出来；它先正常生成文字推理 $t_1$。
  \item 一旦检测到代码结束标记 \verb|</code>|，\textbf{生成就暂停}，把其中的代码 $c_1$ 解析出来、送进代码沙箱去跑。
  \item 沙箱跑完后，把输出 $f_1$（成功结果\emph{或}报错信息）填进 \verb|<interpreter></interpreter>| 标签，回灌给模型。
  \item 模型接着往下生成，直到给出最终答案 $o$、或者再写一段新代码。最终形成一条混合推理轨迹
  $[t_1\oplus c_1\oplus f_1\oplus\cdots\oplus o]$（$\oplus$ 表示拼接）。
\end{enumerate}
\textbf{划重点：} 报错信息也会原样回灌给模型。正是这个“动态反馈”机制，让模型能\emph{从错误中学习}、迭代修正代码——
这是后面“代码自我纠错”涌现行为的来源。

\paragraph{2.3.2 训练细节（工程上的几个关键点）。}
\begin{itemize}[leftmargin=1.6em,itemsep=1pt]
  \item \textbf{框架与超参}：用 VeRL 框架；PPO；冷启动数据训 2 个 epoch；AdamW 优化器，初始学习率 $1\text{e-}6$；
  最大序列长度 16384 token；mini-batch size 512；\textbf{KL 系数设为 0.0}（即不加 KL 约束，让模型放开探索）；
  主基座 Qwen2.5-32B-Instruct。
  \item \textbf{解释器反馈掩码}（Interpreter Feedback Mask）：把 \verb|<interpreter>| 里的反馈内容\emph{从损失计算中屏蔽掉}。
  因为这些是\emph{外部塞进来}的 token、不是模型自己生成的，算进 loss 会干扰训练、破坏模型自己的连贯推理。
  这是个很重要的稳定性技巧。
  \item \textbf{KV-Cache 复用}：每次触发 \verb|</code>| 时，把代码执行前的 KV-cache 全缓存下来，之后只计算并追加
  解释器反馈那段，大幅降低显存开销。
  \item \textbf{异步沙箱}：沙箱以“worker 池”方式运行，各 worker 按自身负载主动拉任务，做到负载均衡、并行执行，
  避免慢线程拖后腿，保证训练吞吐稳定。
\end{itemize}

\subsection{第 3 节 实验：见后文“实验解读”}
（这一节的数据集、基线、主结果与“认知行为分析”，我们集中放在第~\ref{sec:exp} 节展开讲。）

\subsection{第 4 节 背景与相关工作}
\begin{itemize}[leftmargin=1.6em,itemsep=2pt]
  \item \textbf{4.1 LLM 推理}：CoT 由 Wei 等人首创，用自然语言一步步描述来增强推理。研究重心已从“训练时扩展”转向
  \textbf{“测试时扩展”}（test-time scaling，即推理时多花算力换更好结果），手段包括逐步偏好优化（Step-DPO）、
  蒙特卡洛树搜索（MCTS）、RL 等。与 CoT 互补的是 \textbf{PoT}（Program-of-Thought）和 PAL，它们把 Python 解释器等
  外部工具接进来简化与验证复杂步骤。
  \item \textbf{4.2 工具集成推理（TIR）}：TIR 最早是为帮 LLM 解决计算密集的数学题而引入的。Wang 等人提出文字推理与
  代码执行\emph{相互验证}的迭代方法；Chen 等人用 SFT 在自建的“带代码 CoT”数据上集成代码执行——但 SFT 受限于特定数据分布，
  \emph{学不到“何时、如何调用工具”这种自适应策略}。有一项同期工作把 RL 用在 1.5B/7B 的 Qwen2.5-Math 上学工具使用，
  但效果还不够好。\textbf{ReTool 把这条路线放大到了 32B}，并显著超过 Qwen-Math-72B-TIR 和 o1-preview。
\end{itemize}

\subsection{第 5 节 结论}
本文提出 ReTool：一个让 LLM 通过\emph{高效使用代码解释器}来\textbf{自我增强数学推理}的 RL 框架。在 AIME2024/2025 上的
实验表明，ReTool 不仅准确率超过传统文字 RL，而且\textbf{收敛所需训练步数大幅减少}。通过精心的数据构造和专门的工具
使用流水线，ReTool 让模型发展出成熟的“计算介入”策略，为更高效、更强大的工具增强推理铺平了道路。

%========================= 方法图解 =========================
\section{方法图解：一张图看懂 ReTool 的训练回路}

\begin{figure}[H]\centering
  \includegraphics[width=0.97\linewidth]{cand_p04_k00.pdf}
  \caption{文字 RL 训练流程与 ReTool RL 训练流程的对比（原文图 2）。上半 (a) 是普通文字 RL；下半 (b) 是 ReTool。
  图例：Text=文字、Code=代码、Interpreter Feedback=解释器反馈、Final Result=最终结果。}
  \label{fig:arch}
\end{figure}

\noindent\textbf{这张图在说什么、数据怎么流动？}
\begin{itemize}[leftmargin=1.6em,itemsep=2pt]
  \item \textbf{上半 (a) 文字 RL（对照组）}：问题（Question）进 \emph{Policy LLM}（策略模型），它一口气吐出一长串\emph{纯文字}
  （图中灰色块 = Text），到末尾产出最终结果（深蓝小块）→ 算 \emph{Reward}（奖励）→ 转成 \emph{Advantage}（优势）回传更新模型。
  \emph{全程没有代码、没有外部执行}。
  \item \textbf{下半 (b) ReTool（本文方法）}：同样从 Question 出发进 Policy LLM，但生成过程是\textbf{“文字块（灰）+ 代码块（绿）+
  解释器反馈块（蓝）”交错}进行的。每当模型写完一段代码（绿色），就\textbf{往下打到 Code Sandbox（代码沙箱）}里真跑一遍，
  沙箱把运行结果（蓝色）\textbf{回灌}给 Policy LLM；模型据此继续往下想、可能再写新代码……如此循环（图中 “$\cdots$” 表示多轮），
  直到产出最终结果（深蓝）→ 同样走 Reward → Advantage 更新模型。
  \item \textbf{一句话抓住本质}：(a) 是“闷头写作文”，(b) 是“边写边查、写一段就跑一段、跑错了改了再写”。
  而“在哪一步写代码、写什么代码”不是人规定的，是模型在 RL 里\emph{靠最终对错的奖励自己摸索出来的}。
\end{itemize}

%========================= 实验解读 =========================
\section{实验解读：这些数字说明了什么}
\label{sec:exp}

\paragraph{评测设置。} 在 \textbf{AIME2024 \& AIME2025}（美国数学邀请赛，高难数学竞赛题）上评测。为了评测稳定，每个评测集
\textbf{重复跑 32 遍取平均}来估计 pass@1（“一次答对”的准确率）。推理超参：temperature 1.0、top-p 0.7。对比的强基线包括
Qwen2.5-Math-72B-Instruct(-TIR)、Sky-T1、DeepSeek-R1-Zero-Qwen-32B、QwQ-32B-Preview、s1-32B、OpenAI o1-preview，
以及最关键的对照——\textbf{不带工具的文字 RL}（Text-based RL, 同样基于 Qwen2.5-32B-Instruct）。

\begin{table}[H]\centering
\caption{主结果（原文表 1，AIME pass@1，单位 \%）。$\spadesuit$ 文字 RL 也包含了文字版冷启动 SFT 以保证公平对比；
$\diamondsuit$ 冷启动模型在推理时同样接入了代码解释器。“--” 表示原文未报告。}
\begin{tabular}{lcc}
\toprule
模型 & AIME2024 & AIME2025 \\
\midrule
\multicolumn{3}{l}{\textit{已有基线（Existing Baselines）}}\\
Qwen2.5-Math-72B-Instruct        & 30.0 & --   \\
Qwen2.5-Math-72B-Instruct-TIR    & 40.0 & --   \\
Sky-T1                           & 43.3 & --   \\
OpenAI o1-preview                & 44.6 & 37.9 \\
DeepSeek-R1-Zero-Qwen-32B        & 47.0 & --   \\
QwQ-32B-Preview                  & 50.0 & 33.5 \\
s1-32B                           & 56.7 & --   \\
\midrule
\multicolumn{3}{l}{\textit{CI-powered RL（本文方法）}}\\
\textbf{ReTool (Qwen2.5-32B-Instruct)}        & \textbf{67.0} & \textbf{49.3} \\
\textbf{ReTool (DeepSeek-R1-Distill-Qwen-32B)} & \textbf{72.5} & \textbf{54.3} \\
\midrule
\multicolumn{3}{l}{\textit{消融（基于 Qwen2.5-32B-Instruct）}}\\
不训练（Base Model）                & 26.7 & --   \\
去掉 CI（文字 RL $\spadesuit$）       & 40.0 & 36.7 \\
去掉 RL（仅冷启动 $\diamondsuit$）    & 40.9 & 34.5 \\
\bottomrule
\end{tabular}
\end{table}

\paragraph{主结果解读。}
\begin{itemize}[leftmargin=1.6em,itemsep=2pt]
  \item \textbf{又准又快}：ReTool(Qwen2.5-32B) 仅 \textbf{400 步}就到 67.0\% / 49.3\%，而文字 RL 跑 1000+ 步才 40.0\% / 36.7\%。
  这说明“把工具使用显式纳入决策过程”不光提升了上限，还\emph{大幅提升了训练效率}。
  \item \textbf{超过强基线}：AIME2024 上比 s1-32B 高 10.3 个百分点；AIME2025 上比 o1-preview 高 11.4 个百分点。换更强基座
  DeepSeek-R1-Distill-Qwen-32B 后进一步涨到 72.5\% / 54.3\%。
  \item \textbf{消融最有信息量}：把 ReTool 拆开看——\emph{光冷启动}（40.9\%）就已追平\emph{文字 RL}（40.0\%）、远超不训练的
  26.7\%；说明\textbf{数据集本身有效}。而真正把成绩从 40\% 拉到 67\% 的，是\textbf{在带工具的 rollout 上做 RL} 这一步。
  即：\emph{冷启动给底子，RL 才是质变}。
\end{itemize}

\paragraph{认知行为分析（第 3.3 节，本文最有意思的部分）。} 作者在 RL 训练的各个 checkpoint 上统计了一系列“代码相关行为”
指标（原文图 3），发现模型在 RL 过程中\emph{自发}变强：
\begin{itemize}[leftmargin=1.6em,itemsep=2pt]
  \item \textbf{回答更短、更高效}：平均回答长度先\emph{骤降}后缓升，最终比 RL 前\textbf{短约 40\%}（从约 10k 降到约 6k token）。
  原因是“用一段代码替换掉冗长的手算过程”，反而省了 token。
  \item \textbf{越来越爱用代码}：含代码的回答比例（code ratio）一路上升，最终\textbf{近 98\%} 的题都会用到代码。
  \item \textbf{代码越写越复杂}：平均代码行数持续上升，训练结束时约为 RL 前的\textbf{5 倍}；测试集“正确代码”总数从 1k 升到 5k。
  \item \textbf{越用越对}：正确回答里的代码通过率始终\emph{接近 100\%}；错误回答里的代码通过率则呈下降趋势——说明
  “代码能不能跑通”确实影响最终对错。
  \item \textbf{调用时机提前}：代码调用时机（代码起始位置 / 回答总长）在训练中\emph{不断提前}，说明模型学会了\emph{该早用就早用}。
  \item \textbf{“顿悟时刻”——代码自我纠错}（原文图 4）：\emph{在没有任何“纠错示范数据”的情况下}，模型自发学会了改 bug。
  典型案例：模型先写的代码因函数 \texttt{greedy()} 未定义而报 \texttt{NameError}；收到报错后，模型反思
  \emph{“Oops, the functions need to be defined in the same scope. Let's correct that.”}（哎呀，函数得定义在同一作用域里，
  我改一下），随后补全定义、生成了能跑的版本。这是 RL 培养出“元认知”的有力证据。
  \item \textbf{代码用途更多元}（图 5）：RL 前代码主要用于“计算 + 验证”；RL 后用途变得\emph{更丰富多样}，体现出
  “自适应选工具”的能力，也增强了泛化性。
  \item \textbf{对比案例}（图 6）：同一道题，纯文字推理靠“费力的手算”、容易算错而答错；ReTool 则用\emph{简洁代码}替代复杂计算，
  既保证算得对，又让模型能把精力放在整体解题策略上。
\end{itemize}

%========================= takeaways =========================
\section{核心 takeaways 与对 AI 应用开发的启示}

\begin{takeaway}
\textbf{1. “会用工具”可以被 RL 教出来，而且不靠人定规则。} ReTool 最核心的贡献是证明了：模型“何时/如何调用代码”
这种策略，完全可以\emph{仅凭最终结果的奖励}自学成才，比人工写规则或纯 SFT 模仿更鲁棒、更能泛化。
\end{takeaway}

\begin{takeaway}
\textbf{2. 把“能验证的活儿”交给工具，把“想清楚”留给模型。} 精确计算、枚举、验证这类事，代码做得又快又准；
模型负责高层推理与决策。这种“神经 + 符号”的分工，对所有\emph{需要精确性}的 AI 应用（数据分析、财务计算、
代码辅助、科学计算）都是可直接借鉴的范式。
\end{takeaway}

\begin{takeaway}
\textbf{3. 奖励越简单越好，警惕 reward hacking。} 作者特意\emph{只奖励最终答案对错}、不给“代码能跑”单独加分，
就是怕模型钻空子（比如为了拿“可执行”分而乱写代码）。做 RLVR 时，\textbf{选一个干净、可自动验证、难以作弊的奖励信号}
往往比堆砌复杂奖励更重要。
\end{takeaway}

\begin{takeaway}
\textbf{4. 工程细节决定成败：反馈掩码 + KV-Cache 复用 + 异步沙箱。} 想自己搭“带真实工具执行的 RL”，这三点是关键：
\emph{把外部塞进来的 token 从 loss 里屏蔽}（否则训练不稳）、\emph{复用代码执行前的 KV-cache}（省显存）、
\emph{异步并行的沙箱池}（别让代码执行拖慢训练）。这些是论文里很实在的可复用经验。
\end{takeaway}

\begin{takeaway}
\textbf{5. 冷启动给底子，RL 给质变。} 消融实验清楚地告诉我们：好的示范数据（冷启动 SFT）能让模型“及格”，
但要从 40\% 跨到 67\%，必须靠\emph{在真实工具交互上做 RL}。做 agent / 工具调用产品时，别指望只靠造数据微调就到顶，
真正拉开差距的是“让模型在带真实反馈的环境里试错学习”。
\end{takeaway}

\begin{takeaway}
\textbf{6. “顿悟时刻”是涌现的。} 代码自我纠错没有被专门教过，却在 RL 中自发出现。这提示我们：只要给对\emph{环境
（能拿到真实报错反馈）+ 奖励（结果导向）}，模型能涌现出超出训练数据的高级行为。把“真实的失败反馈”喂回给模型，
是培养这种能力的关键。
\end{takeaway}

%========================= 参考文献 =========================
\section*{参考文献（仅列最关键几条，英文原样）}
\begin{enumerate}[leftmargin=2em,itemsep=1pt]
\item DeepSeek-AI et al. \emph{DeepSeek-R1: Incentivizing reasoning capability in LLMs via reinforcement learning}, 2025. \url{https://arxiv.org/abs/2501.12948}
\item John Schulman, Filip Wolski, Prafulla Dhariwal, Alec Radford, Oleg Klimov. \emph{Proximal Policy Optimization Algorithms}, 2017. \url{https://arxiv.org/abs/1707.06347}
\item Wenhu Chen, Xueguang Ma, Xinyi Wang, William W. Cohen. \emph{Program of Thoughts Prompting: Disentangling computation from reasoning for numerical reasoning tasks}, 2023. \url{https://arxiv.org/abs/2211.12588}
\item Xuefeng Li, Haoyang Zou, Pengfei Liu. \emph{ToRL: Scaling tool-integrated RL}, 2025. \url{https://arxiv.org/abs/2503.23383}
\item Qwen et al. \emph{Qwen2.5 Technical Report}, 2025. \url{https://arxiv.org/abs/2412.15115}
\item OpenAI. \emph{Learning to reason with LLMs}, 2024. \url{https://openai.com/index/learning-to-reason-with-llms/}
\end{enumerate}

\end{document}
